% =========================================================================
% Deadline-Aware Multi-Agent Deep Reinforcement Learning for
% Cooperative Edge Caching under Partial Observability
%
% IEEE Conference Paper — IEEEtran class
% =========================================================================
\documentclass[conference]{IEEEtran}

\usepackage{amsmath,amssymb}
\usepackage{booktabs}
\usepackage{graphicx}
\usepackage{xcolor}
\usepackage{algorithm}
\usepackage{algpseudocode}
\usepackage{hyperref}
\usepackage{cite}
\usepackage{array}
\usepackage{multirow}
\setlength{\emergencystretch}{3em}

\hypersetup{
  colorlinks=true,
  linkcolor=blue,
  citecolor=blue,
  urlcolor=blue
}

% =========================================================================
\title{Deadline-Aware Multi-Agent Deep Reinforcement Learning for
       Cooperative Edge Caching under Partial Observability}

\author{
  \IEEEauthorblockN{Sanskar Dhyani, Soniya, Shiv Sunder Behera}
  \IEEEauthorblockA{Department of Computer Science and Engineering \\
    Dr.\ B.\ R.\ Ambedkar National Institute of Technology Jalandhar \\
    Jalandhar, Punjab, India}
  \\[4pt]
  \IEEEauthorblockN{\textit{Supervisor: Dr.\ Someshula Manoj Kumar}}
  \IEEEauthorblockA{Department of Computer Science and Engineering \\
    Dr.\ B.\ R.\ Ambedkar National Institute of Technology Jalandhar \\
    Jalandhar, Punjab, India}
}

\begin{document}
\maketitle

% =========================================================================
\begin{abstract}
Edge computing networks must serve heterogeneous content workloads
under strict latency deadlines while operating with only partial
knowledge of the global system state.  Existing deep reinforcement
learning (DRL) caching approaches, such as the Intelligent Caching at
the Edge (ICE) framework, assume a single centralised agent with full
Markov-state visibility and static Zipf popularity—assumptions that do
not hold in realistic distributed deployments.  We present
\emph{DEADLINE\_DCCC}, a Deadline-Aware Cooperative Caching framework
that replaces these assumptions with (i)~a Partially Observable Markov
Decision Process (POMDP) formulation in which each edge node maintains
an independent DQN agent acting on a seven-dimensional local
observation, (ii)~a Newton-cooling popularity model that captures both
historical and recency-weighted demand, and (iii)~a six-term
deadline-aware scoring function that jointly optimises cache locality,
load balance, content replication overlap, path latency, and deadline
risk.  Evaluation on a heterogeneous five-node cluster (capacities
$[9,9,2.7,9,2.7]$\,GB) over 6\,000 requests with region-skewed Zipf
traffic achieves a 54.5\% hit rate, a 54.5\% deadline satisfaction
rate, and an average end-to-end latency of 87.06\,ms, with
P95\,=\,175.4\,ms.
\end{abstract}

\begin{IEEEkeywords}
edge caching, deep reinforcement learning, cooperative caching,
partial observability, POMDP, deadline-aware routing, multi-agent
systems, content delivery networks
\end{IEEEkeywords}

% =========================================================================
\section{Introduction}
\label{sec:intro}

The proliferation of latency-sensitive applications—live video
streaming, augmented reality, industrial IoT, and remote healthcare—
has driven substantial research interest in \emph{mobile edge
computing} (MEC), where computation and storage resources are pushed to
the network periphery~\cite{shi2016edge}.  Content Delivery Networks
(CDNs) increasingly rely on cooperative edge caches to reduce
backbone traffic and guarantee quality-of-service (QoS) targets
expressed as per-request latency deadlines.

Traditional caching policies such as Least Recently Used (LRU) and
Least Frequently Used (LFU) are agnostic to content popularity
dynamics and network topology.  Machine learning approaches, and
particularly deep reinforcement learning (DRL), have been applied to
learn adaptive caching policies directly from interactions with the
system~\cite{sadeghi2019optimal,zhong2020deep}.  The ICE framework
~\cite{ice2022} proposed a single DQN agent operating on a global MDP
state to jointly control content placement and eviction, demonstrating
significant improvements in hit rate and energy consumption over static
baselines.

However, ICE and similar single-agent approaches suffer from three
structural limitations when deployed in realistic distributed edge
networks.  First, \emph{centralised state} is expensive to maintain:
aggregating cache inventories, load levels, and latency measurements
from all nodes requires constant inter-node signalling that itself
consumes bandwidth and introduces staleness~\cite{liu2019decentralised}.
Second, \emph{static popularity models} (Zipf with fixed exponent) fail
to track content lifecycle: viral items spike rapidly and cool off,
whereas long-tail items maintain baseline demand for extended
periods~\cite{traverso2013temporal}.  Third, existing formulations
ignore \emph{per-request deadlines}, which are increasingly mandated by
service-level agreements for real-time applications.

This paper presents \emph{DEADLINE\_DCCC} (Deadline-aware
Decentralised Cooperative Caching with DRL), which addresses all three
limitations.  Our contributions are:

\begin{enumerate}
  \item A POMDP formulation in which each edge node $i$ maintains an
    independent DQN agent observing only a seven-dimensional local
    state vector, eliminating the need for global state
    synchronisation.
  \item A Newton-cooling popularity model that blends Zipf-prior
    popularity with a time-decaying request count, providing a
    dynamically recency-weighted demand signal.
  \item A six-term deadline-aware scoring function that simultaneously
    accounts for cache locality, content overlap, path latency,
    quadratic load cost, and residual deadline risk.
  \item A two-tier routing architecture ($\mathrm{edge} \to
    \mathrm{cooperative} \to \mathrm{cloud}$) with strict
    ($D_1 = 35$\,ms) and cooperative ($D_2 = 95$\,ms) deadline tiers.
  \item An empirical evaluation over heterogeneous node capacities and
    region-skewed traffic showing competitive hit rates alongside
    explicit deadline satisfaction reporting.
\end{enumerate}

The remainder of the paper is organised as follows.  Section~\ref{sec:related}
reviews related work.  Section~\ref{sec:model} describes the system
model.  Section~\ref{sec:pomdp} provides the POMDP formulation.
Section~\ref{sec:algorithm} presents the DEADLINE\_DCCC algorithm.
Section~\ref{sec:differences} explicitly compares our approach to the
ICE base paper.  Section~\ref{sec:eval} reports experimental results.
Section~\ref{sec:conclusion} concludes.

% =========================================================================
\section{Related Work}
\label{sec:related}

\subsection{Adaptive Edge Caching with DRL}

Sadeghi et al.~\cite{sadeghi2019optimal} showed that model-free DRL can
learn near-optimal caching policies without explicit knowledge of
content popularity distributions, outperforming classical replacement
rules on synthetic and real traces.  Zhong et al.~\cite{zhong2020deep}
applied DQN to proactive caching in heterogeneous networks, exploiting
temporal patterns in user requests to pre-populate caches before demand
peaks.

The ICE framework~\cite{ice2022} unified placement, eviction, and
routing decisions under a single DQN trained on an MDP with global
edge-state observations.  ICE introduced the \emph{ICE reward}
(popularity $\times$ $\log(1 + \mathrm{size})$) as an eviction metric
that favours retaining large popular items and demonstrated that
DRL-based policies outperform LRU and LFU on hit rate and energy across
a range of Zipf exponents.  Our work builds directly on ICE's reward
formulation but extends it to the multi-agent, deadline-constrained
setting.

\subsection{Cooperative Caching}

Content replication across cooperative edge nodes can multiply
effective cache capacity, but naive replication wastes space and
saturates inter-node links~\cite{borst2010distributed}.  Femtocaching
~\cite{golrezaei2013femtocaching} showed analytically that cooperative
caching across small-cell base stations can dramatically reduce
backhaul load.  More recent work has applied multi-agent RL to
cooperative caching~\cite{zheng2018heterogeneous}, but typically
assumes homogeneous node capacity and ignores the partial observability
problem.  We explicitly model heterogeneous capacity
($[9,9,2.7,9,2.7]$\,GB) and penalise redundant replication through the
overlap term in our scoring function.

\subsection{Partial Observability in Multi-Agent RL}

The Partially Observable Markov Decision Process (POMDP) provides the
correct theoretical framework when agents cannot observe the full system
state~\cite{kaelbling1998planning}.  Lowe et al.~\cite{lowe2017multi}
introduced the \emph{centralised training, decentralised execution}
(CTDE) paradigm; however, their approach still requires sharing agent
observations at training time.  Our system avoids inter-agent
communication entirely during both training and execution, making it
suitable for deployments where nodes may belong to different
administrative domains.

\subsection{Deadline-Aware Content Delivery}

Meeting per-request latency deadlines has been studied in the
real-time scheduling literature~\cite{buttazzo2011hard}, but its
integration with learning-based caching is limited.  Xu
et al.~\cite{xu2019joint} considered joint offloading and caching with
delay constraints in MEC, while Liu et al.~\cite{liu2020deadline}
proposed deadline-constrained task scheduling on edge servers.  Neither
work addresses content popularity dynamics or multi-tier cooperative
routing.  DEADLINE\_DCCC closes this gap by embedding deadline risk
directly into the node-selection score and the DRL reward.

% =========================================================================
\section{System Model}
\label{sec:model}

\subsection{Network Topology}

We consider a three-tier content delivery infrastructure: a cloud
origin server, an edge cluster of $N = 5$ heterogeneous edge nodes
$\mathcal{N} = \{n_0, n_1, n_2, n_3, n_4\}$, and a population of
users organised into $R = 3$ geographic regions.  Edge nodes are
connected to each other via a low-latency intra-cluster fabric and to
the cloud via a higher-latency backhaul link.  Nodes $n_0, n_1, n_3$
each have a storage capacity of 9\,GB, while $n_2$ and $n_4$ have
2.7\,GB each, reflecting the heterogeneous hardware typical of real
edge deployments.

\subsection{Latency Model}

For a request from region $u$ served locally by node $n$ holding
content $c$ of size $s_c$\,MB, the edge delivery latency is

\begin{equation}
  \ell_{\mathrm{edge}}(n,c,u) =
    \ell_{\mathrm{base}}
    + \alpha_{\mathrm{tx}} \cdot s_c
    + \alpha_{\mathrm{load}} \cdot \mathrm{load}(n)
    + \alpha_{\mathrm{region}} \cdot |r(n) - u|,
  \label{eq:latency}
\end{equation}

where $\ell_{\mathrm{base}} = 8$\,ms is the base propagation and
scheduling cost, $\alpha_{\mathrm{tx}} = 0.012$\,ms/MB models a
1\,Gbps edge link, $\alpha_{\mathrm{load}} = 10$\,ms is a queuing
penalty, and $\alpha_{\mathrm{region}} = 4$\,ms/hop penalises
cross-region traversal.  The load $\mathrm{load}(n) \in [0, 1.5]$ is
a continuous variable updated after each request.

Cooperative serving from a peer node $n'$ incurs an additional fixed
cluster-traversal cost $\ell_{\mathrm{coop}} = 35$\,ms plus half the
local edge latency at the holder:
$\ell_{\mathrm{cooperative}}(n') = \ell_{\mathrm{coop}} +
\frac{1}{2}\ell_{\mathrm{edge}}(n',c,u)$.

Cloud fetching costs $\ell_{\mathrm{cloud}}(c) = 120 + 0.025 \cdot s_c$
ms, reflecting a slower 40\,Mbps backhaul link.

\subsection{Load Dynamics}

Each non-cloud serve increments the serving node's load by 0.25.
After every request, loads are decayed by a factor of 0.9 and a small
uniform noise $\epsilon \sim \mathcal{U}(-0.005, 0.005)$ is added to
model background variability:

\begin{equation}
  \mathrm{load}(n, t+1) = 0.9 \cdot \mathrm{load}(n, t) + 0.25 \cdot
  \mathbf{1}[\text{served by } n] + \epsilon.
  \label{eq:load}
\end{equation}

\subsection{Content Model and Newton-Cooling Popularity}

The content catalogue consists of $M = 120$ items drawn from four
types: video, image, IoT, and audio.  Content sizes are uniformly
distributed in $[50, 2500]$\,MB.  Each item $c$ has a base Zipf
popularity $\rho_c$ with exponent $\alpha_z = 0.85$ and a region
affinity bias in $[0.1, 1.0]$.

To capture the lifecycle of content demand, we introduce a
\emph{Newton-cooling popularity} model that blends the Zipf prior with
a decaying request-count signal:

\begin{equation}
  p_c(t) = \left(\rho_c + \frac{q_c(t)}{20}\right) e^{-\lambda t},
  \label{eq:newton}
\end{equation}

where $q_c(t)$ is the cumulative request count for content $c$ up to
time $t$ and $\lambda = 4 \times 10^{-4}$ is the cooling coefficient.
This formulation ensures that (i)~recently requested content benefits
from a positive $q_c$ boost, and (ii)~all content eventually cools,
preventing the cache from permanently favouring stale popular items.
The ICE-style eviction metric $\rho_c \cdot \log(1 + s_c)$ is then
evaluated using $p_c(t)$ in place of the static popularity.

\subsection{Traffic Model}

Requests arrive according to a Zipf distribution over contents
conditioned on the user's region.  Region weights are
$[0.60, 0.25, 0.15]$, modelling a hot primary region, a moderate
secondary region, and a light tertiary region.  A cross-region
probability of 0.35 models users who occasionally request content
hosted in a different region.  The request stream is seeded for
reproducibility; all policies are evaluated against an identical
stream.

\subsection{Deadline Tiers}

Two latency deadline tiers are defined:
\begin{itemize}
  \item \textbf{Strict edge deadline} $D_1 = 35$\,ms: content that
    can be served locally within this budget.
  \item \textbf{Cooperative deadline} $D_2 = 95$\,ms: content that
    can be served via a cooperative peer within this budget.
\end{itemize}
Any request whose end-to-end latency exceeds $D_2$ triggers a cloud
fetch, which is never considered deadline-satisfying.

% =========================================================================
\section{Problem Formulation}
\label{sec:pomdp}

\subsection{Limitations of the MDP Assumption}

The ICE framework~\cite{ice2022} models the caching problem as a
standard Markov Decision Process, where a single agent observes the
full global state including all node caches, load levels, and
content popularity counters.  While this simplifies training, the MDP
assumption breaks down in practice because (i)~global state collection
incurs $O(N \cdot M)$ inter-node communication per request, (ii)~the
centralised agent becomes a single point of failure, and (iii)~the
state space scales prohibitively with cluster size.

\subsection{POMDP Formulation}

We model the cooperative caching problem as a \emph{Partially
Observable Markov Decision Process} (POMDP), defined by the tuple
$\langle \mathcal{S}, \mathcal{A}, \mathcal{T}, \mathcal{R},
\mathcal{O}, \mathcal{Z}, \gamma \rangle$:

\begin{itemize}
  \item \textbf{State space} $\mathcal{S}$: The full system state at
    time $t$ comprises the cache contents of all nodes, their load
    vectors, and the current popularity estimates for all content items.
    This state is \emph{not} fully observed by any single agent.

  \item \textbf{Observation space} $\mathcal{O}$: Agent $i$ at node
    $n_i$ observes the seven-dimensional local observation vector
    $\mathbf{o}_i(t) \in [0,1]^7$ defined as:
    \begin{equation}
      \mathbf{o}_i(t) = \Bigl[
        \rho_c,\;
        \tfrac{s_c}{2500},\;
        h_i,\;
        \mathrm{load}_i,\;
        \tfrac{\ell_i}{200},\;
        \Omega_i,\;
        \tfrac{\max(0,\ell_i - D_1)}{200}
      \Bigr],
      \label{eq:obs}
    \end{equation}
    where $\rho_c$ is the base Zipf popularity of the requested content,
    $s_c$ is its size in MB normalised by the maximum (2500\,MB),
    $h_i \in \{0,1\}$ indicates whether node $n_i$ currently holds
    the content, $\mathrm{load}_i$ is the node's current load,
    $\ell_i$ is the estimated edge delivery latency normalised by
    200\,ms, $\Omega_i$ is the cache overlap ratio (defined in
    Section~\ref{sec:algorithm}), and the final term is the deadline
    risk normalised by 200\,ms.

  \item \textbf{Action space} $\mathcal{A}$: Each agent selects from
    three routing actions:
    \begin{equation}
      \mathcal{A} = \{0: \textsc{serve\_local},\;
                       1: \textsc{forward\_dccc},\;
                       2: \textsc{default\_policy}\}.
    \end{equation}
    \textsc{serve\_local} routes the request to the local node.
    \textsc{forward\_dccc} triggers cooperative peer scoring.
    \textsc{default\_policy} delegates to the deterministic
    deadline-aware scoring heuristic.

  \item \textbf{Transition model} $\mathcal{T}$: Transitions are
    \emph{partially known} to each agent.  The agent can estimate its
    own cache state, load, and latency, but cannot observe changes to
    peer nodes' caches induced by other requests handled elsewhere in
    the cluster.

  \item \textbf{Reward function} $\mathcal{R}$: The scalar reward
    received by the agent responsible for routing request $r$ is:
    \begin{equation}
      \mathcal{R}(s, a) =
        r_{\mathrm{hit}}
        - 0.03 \cdot \ell_r
        + 5 \cdot \mathbf{1}[\ell_r \leq D_{\mathrm{tier}}]
        - 5 \cdot \mathbf{1}[\ell_r > D_{\mathrm{tier}}],
      \label{eq:reward}
    \end{equation}
    where $r_{\mathrm{hit}} \in \{+5, +3, -5\}$ for edge hit,
    cooperative hit, and cloud fetch respectively, $\ell_r$ is the
    realised latency in milliseconds, and $D_{\mathrm{tier}}$ is the
    applicable deadline ($D_1$ or $D_2$).

  \item \textbf{Discount factor} $\gamma = 0.95$: Balances immediate
    latency penalties against long-term cache placement benefits.
\end{itemize}

\subsection{Decoupled Multi-Agent Training}

Five independent DQN agents are trained simultaneously, one per node.
Each agent observes only its own local state and receives its own
scalar reward.  Agents do not share parameters or communicate during
training.  This constitutes a \emph{decentralised execution and
decentralised training} (DEDT) paradigm, which is more restrictive than
the CTDE paradigm of MADDPG~\cite{lowe2017multi} but avoids the need
for an inter-agent communication channel entirely.  The independence
assumption is justified because each node's routing decision
predominantly depends on its own cache state and load, with only a weak
coupling through cooperative serving paths.

% =========================================================================
\section{Proposed Algorithm: DEADLINE\_DCCC}
\label{sec:algorithm}

\subsection{Six-Term Deadline-Aware Scoring Function}

The core contribution of this paper is the \emph{deadline-aware
scoring function} used to select the best edge node for each request.
For node $n$, content $c$, and user region $u$, the score is:

\begin{align}
  \sigma(n, c, u) &=
    \underbrace{2.0 \cdot p_c(t)}_{\text{popularity}}
    + \underbrace{6.0 \cdot \chi(n, c)}_{\text{cache hit}}
    - \underbrace{0.5 \cdot \Omega_n}_{\text{overlap}} \nonumber \\
    &\quad
    - \underbrace{0.08 \cdot \ell(n,c,u)}_{\text{latency}}
    - \underbrace{1.2 \cdot \mathrm{load}(n)^2}_{\text{load}}
    - \underbrace{0.4 \cdot \delta(n,c,u)}_{\text{deadline risk}},
  \label{eq:score}
\end{align}

where:
\begin{itemize}
  \item $p_c(t)$ is the Newton-cooling popularity (Eq.~\ref{eq:newton}).
  \item $\chi(n,c) = 1.0$ if $n$ currently holds content $c$, otherwise
    $\min(0.5, p_c(t))$.  This term strongly rewards routing to a node
    that already caches the content while still providing a
    popularity-proportional incentive for nodes without it.
  \item $\Omega_n = |\mathcal{C}_n \cap \mathcal{C}_{-n}| /
    \max(1, |\mathcal{C}_n|)$ is the \emph{cache overlap ratio}: the
    fraction of node $n$'s cache contents that are also held by at
    least one other cluster node.  Penalising overlap encourages
    \emph{soft partitioning}, increasing effective cluster capacity.
  \item $\ell(n, c, u)$ is the estimated edge delivery latency
    (Eq.~\ref{eq:latency}).
  \item $\mathrm{load}(n) \in [0, 1.5]$ is the node's current load.
    The quadratic term $\mathrm{load}^2$ applies a negligible penalty
    at moderate utilisation ($\mathrm{load}=0.5 \Rightarrow -0.3$
    points) but a strong penalty at saturation ($\mathrm{load}=1.5
    \Rightarrow -2.7$ points), enabling the policy to route around
    overloaded nodes without abandoning lightly loaded ones.
  \item $\delta(n,c,u) = \max(0, \ell(n,c,u) - D_1)$ is the residual
    deadline risk, only active when the estimated latency exceeds the
    strict deadline $D_1 = 35$\,ms.
\end{itemize}

The coefficient magnitudes were calibrated on a single seed (seed 42)
and held fixed across all evaluation seeds to prevent overfitting to
the training distribution.  An ablation variant \textsc{deadline\_no\_overlap}
removes the $\Omega_n$ term to isolate the contribution of the
overlap penalty.

\subsection{Two-Tier Routing Protocol}

Algorithm~\ref{alg:routing} describes the complete DEADLINE\_DCCC
routing decision for an incoming request.

\begin{algorithm}[t]
\caption{DEADLINE\_DCCC Routing Decision}
\label{alg:routing}
\begin{algorithmic}[1]
\Require request $(c, u)$, edge cluster $\mathcal{N}$, deadlines $D_1, D_2$,
         DQN agents $\{\pi_i\}_{i \in \mathcal{N}}$, time step $t$
\Ensure  served latency $\ell$, hit type $h$

\State $n_{\mathrm{local}} \leftarrow$ home node for region $u$
\State $\mathbf{o} \leftarrow$ \Call{BuildObservation}{$n_{\mathrm{local}}$, $c$, $t$}
\State $a \leftarrow \pi_{n_{\mathrm{local}}}(\mathbf{o})$ \Comment{DQN action}

\If{$a = \textsc{serve\_local}$ \textbf{or} $n_{\mathrm{local}}.\mathrm{has}(c)$}
  \State $\ell \leftarrow \ell_{\mathrm{edge}}(n_{\mathrm{local}}, c, u)$
  \If{$\ell \leq D_1$}
    \State \Return $\ell$, \textsc{edge\_hit}
  \EndIf
\EndIf

\If{$a = \textsc{forward\_dccc}$ \textbf{or} $a = \textsc{default\_policy}$}
  \State $p_c \leftarrow$ \Call{NewtonPopularity}{$c$, $t$}
  \For{each node $n \in \mathcal{N}$}
    \State $\Omega_n \leftarrow$ \Call{CacheOverlap}{$n$, $\mathcal{N}$}
    \State $\ell_n \leftarrow \ell_{\mathrm{edge}}(n, c, u)$
    \State $\sigma_n \leftarrow$ \Call{DeadlineScore}{$n$, $c$, $p_c$, $\ell_n$, $\Omega_n$, $D_1$}
  \EndFor
  \State $n^* \leftarrow \arg\max_{n \in \mathcal{N}} \sigma_n$
  \State $\ell \leftarrow \ell_{\mathrm{cooperative}}(n^*)$
  \If{$\ell \leq D_2$}
    \If{$n^*.\mathrm{has}(c)$}
      \State \Return $\ell$, \textsc{coop\_hit}
    \Else
      \State \Call{ICEEvict}{$n^*$, $c$, $p_c$} \Comment{Cache content at best node}
      \State \Return $\ell$, \textsc{coop\_hit}
    \EndIf
  \EndIf
\EndIf

\State $\ell \leftarrow \ell_{\mathrm{cloud}}(c)$
\State \Return $\ell$, \textsc{cloud\_fetch}

\Statex
\Function{DeadlineScore}{$n$, $c$, $p_c$, $\ell_n$, $\Omega_n$, $D_1$}
  \State $\chi \leftarrow 1.0$ if $n.\mathrm{has}(c)$ else $\min(0.5, p_c)$
  \State $\delta \leftarrow \max(0, \ell_n - D_1)$
  \State \Return $2.0 p_c + 6.0 \chi - 0.5 \Omega_n - 0.08 \ell_n
                  - 1.2 \cdot n.\mathrm{load}^2 - 0.4 \delta$
\EndFunction
\Statex
\Function{ICEEvict}{$n$, $c$, $p_c$}
  \If{$n.\mathrm{cache}$ is full}
    \State $c' \leftarrow \arg\min_{x \in n.\mathrm{cache}}\; p_x(t) \cdot
           \log(1 + s_x)$
    \State Evict $c'$ from $n$
  \EndIf
  \State Insert $c$ into $n.\mathrm{cache}$
\EndFunction
\end{algorithmic}
\end{algorithm}

\subsection{Independent DQN Agents}

Each node $n_i$ maintains a DQN with the architecture
$7 \to 64 \to \mathrm{ReLU} \to 64 \to \mathrm{ReLU} \to 3$.  The
policy network $Q(s, a;\,\theta_i)$ and a periodically synchronised
target network $Q(s, a;\,\theta_i^-)$ are used to compute the
Bellman target:

\begin{equation}
  y_j = r_j + \gamma \max_{a'} Q(s'_j, a';\,\theta_i^-),
  \label{eq:bellman}
\end{equation}

and the MSE loss $\mathcal{L}(\theta_i) = \mathbb{E}\bigl[(y_j -
Q(s_j, a_j;\,\theta_i))^2\bigr]$ is minimised by Adam with learning
rate $\eta = 10^{-3}$.

Exploration follows an $\varepsilon$-greedy policy with
$\varepsilon_0 = 1.0$, $\varepsilon_{\min} = 0.05$, and per-episode
decay factor $\kappa = 0.97$.  Decaying $\varepsilon$ once per episode
(rather than per step) prevents premature collapse of exploration
within the first training episode.  The target network is synchronised
with the policy network once per episode.  Training runs for 50
episodes of 1\,000 requests each; evaluation uses the full 6\,000-request
stream with $\varepsilon = 0$.

\subsection{ICE-Style Eviction}

Cache eviction inherits the ICE-style utility metric:
\begin{equation}
  \nu(c, t) = p_c(t) \cdot \log(1 + s_c),
  \label{eq:ice_evict}
\end{equation}
where $p_c(t)$ is the Newton-cooling popularity and $s_c$ is the
content size in MB.  When node $n^*$ must make room for newly fetched
content, the item with the lowest $\nu$ is evicted.  This metric
favours retaining large popular items over small unpopular ones,
which is appropriate given the 50--2500\,MB content size range.

% =========================================================================
\section{Differences from the ICE Base Paper}
\label{sec:differences}

Table~\ref{tab:comparison} summarises the principal design differences
between DEADLINE\_DCCC and the ICE base paper~\cite{ice2022}.

\begin{table}[t]
  \caption{Design Comparison: ICE vs.\ DEADLINE\_DCCC}
  \label{tab:comparison}
  \renewcommand{\arraystretch}{1.25}
  \begin{tabular}{lp{1.8cm}p{2.4cm}}
    \toprule
    \textbf{Aspect} & \textbf{ICE~\cite{ice2022}} & \textbf{DEADLINE\_DCCC (Ours)} \\
    \midrule
    Observability    & MDP (global state)          & POMDP (local state per node) \\
    Agents           & Single centralised DQN      & 5 independent per-node DQNs \\
    Popularity model & Static Zipf                 & Newton-cooling: $({\rho_c + q_c/20}) e^{-\lambda t}$ \\
    Eviction metric  & $\rho_c \cdot \log(1{+}s_c)$ & Same, with dynamic $p_c(t)$ \\
    Node selection   & Nearest/hash                & 6-term deadline-aware score \\
    Deadlines        & None                        & $D_1=35$\,ms, $D_2=95$\,ms \\
    Routing tiers    & Single-tier edge            & Edge $\to$ coop $\to$ cloud \\
    Capacity         & Homogeneous                 & $[9,9,2.7,9,2.7]$\,GB \\
    Traffic          & Uniform Zipf                & Region-skewed: $[0.60,0.25,0.15]$ \\
    Metrics          & Hit rate, energy            & + deadline satisfaction, P95/P99, load imbalance \\
    \bottomrule
  \end{tabular}
\end{table}

Several design choices merit further discussion.

\textbf{POMDP vs.\ MDP.}  Replacing the single centralised MDP agent
with per-node POMDP agents eliminates the $O(N \cdot M)$ state
broadcast cost.  Each agent's seven-dimensional observation
(Eq.~\ref{eq:obs}) can be computed entirely from local measurements,
making the system scalable to large clusters.

\textbf{Newton-cooling popularity.}  The static Zipf model in ICE
assigns a fixed demand rank to each content item for the lifetime of
the simulation.  In practice, content demand follows a lifecycle:
items become popular quickly and cool off over time.
Equation~\ref{eq:newton} models this by introducing a time-decaying
exponential multiplier.  The cooling coefficient $\lambda = 4 \times
10^{-4}$ was chosen so that a content item's popularity halves in
approximately 1\,733 request ticks, matching the observed decay
pattern in the request stream.

\textbf{Overlap penalty.}  Without a mechanism to discourage
replication, the ICE serve-then-cache rule tends toward full
replication of the top-$k$ items across all nodes, which wastes the
heterogeneous capacity of the weaker (2.7\,GB) nodes.  The overlap
term $-0.5 \Omega_n$ in Eq.~\ref{eq:score} penalises routing to nodes
whose caches are already largely replicated elsewhere, nudging the
cluster toward soft content partitioning.

\textbf{Quadratic load penalty.}  The $-1.2 \cdot \mathrm{load}^2$
term applies a saturation-only penalty: at $\mathrm{load} = 0.5$
the contribution is only $-0.3$ score points, while at
$\mathrm{load} = 1.5$ it reaches $-2.7$ points.  This avoids
prematurely routing away from a node that is lightly loaded (where
the round-trip latency would still be good) while strongly penalising
fully saturated nodes.

\textbf{Deadline risk as a scoring term.}  Rather than implementing
deadlines as hard constraints (which can lead to infeasibility under
load), we embed deadline risk $\delta = \max(0, \ell - D_1)$ as a
soft penalty term.  This gracefully degrades under peak load: when no
node can meet $D_1$, the system selects the node with the smallest
latency excess rather than failing.

% =========================================================================
\section{Experimental Evaluation}
\label{sec:eval}

\subsection{Experimental Setup}

All experiments are conducted using a discrete-event simulator
implemented in Python.  The simulator models atomic request processing
with the latency model of Eq.~\ref{eq:latency} and the load dynamics
of Eq.~\ref{eq:load}.  Key parameters are summarised in
Table~\ref{tab:params}.  DRL agents are implemented using PyTorch with
CPU-only execution for reproducibility (GPU floating-point
non-determinism is avoided).

\begin{table}[t]
  \caption{Simulation Parameters}
  \label{tab:params}
  \renewcommand{\arraystretch}{1.2}
  \begin{tabular}{lll}
    \toprule
    \textbf{Parameter} & \textbf{Symbol} & \textbf{Value} \\
    \midrule
    Edge nodes                  & $N$          & 5 \\
    Node capacities (GB)        & --           & $[9,9,2.7,9,2.7]$ \\
    Content catalogue size      & $M$          & 120 \\
    Content size range (MB)     & $[s_{\min}, s_{\max}]$ & $[50, 2500]$ \\
    Zipf exponent               & $\alpha_z$   & 0.85 \\
    Region weights              & --           & $[0.60, 0.25, 0.15]$ \\
    Cross-region probability    & --           & 0.35 \\
    Strict deadline             & $D_1$        & 35\,ms \\
    Cooperative deadline        & $D_2$        & 95\,ms \\
    Base edge latency           & $\ell_{\mathrm{base}}$ & 8\,ms \\
    Coop. path overhead         & $\ell_{\mathrm{coop}}$ & 35\,ms \\
    Cloud base latency          & --           & 120\,ms \\
    Cooling coefficient         & $\lambda$    & $4 \times 10^{-4}$ \\
    Total evaluation requests   & --           & 6\,000 \\
    DRL training episodes       & --           & 50 \\
    Requests per episode        & --           & 1\,000 \\
    DRL learning rate           & $\eta$       & $10^{-3}$ \\
    Discount factor             & $\gamma$     & 0.95 \\
    Initial $\varepsilon$       & $\varepsilon_0$ & 1.0 \\
    Minimum $\varepsilon$       & $\varepsilon_{\min}$ & 0.05 \\
    $\varepsilon$ decay (per episode) & $\kappa$ & 0.97 \\
    Batch size                  & --           & 64 \\
    Hidden layer width          & --           & 64 \\
    \bottomrule
  \end{tabular}
\end{table}

We compare DEADLINE\_DCCC against the following baselines:

\begin{itemize}
  \item \textbf{LRU}: Least Recently Used eviction with hash-based
    node selection.
  \item \textbf{LFU}: Least Frequently Used eviction with hash-based
    node selection.
  \item \textbf{DCCC\_HASH}: DCCC cooperative caching with consistent
    hash routing and ICE-style eviction, without deadline awareness.
  \item \textbf{DEADLINE\_HEURISTIC}: The six-term scoring function
    (Eq.~\ref{eq:score}) without DRL augmentation.
  \item \textbf{DEADLINE\_NO\_OVERLAP}: Ablation that removes the
    overlap term ($\Omega_n = 0$ always) to isolate its contribution.
\end{itemize}

\subsection{Performance Results}

Table~\ref{tab:results} reports the aggregate evaluation metrics over
6\,000 requests.

\begin{table}[t]
  \caption{Policy Performance Comparison (6\,000 requests, seed 42)}
  \label{tab:results}
  \renewcommand{\arraystretch}{1.2}
  \begin{tabular}{lcccc}
    \toprule
    \textbf{Policy} & \textbf{Hit Rate} & \textbf{Deadline} & \textbf{Avg.\ Lat.} & \textbf{Cloud} \\
    & (\%) & \textbf{Sat.\ (\%)} & \textbf{(ms)} & \textbf{(\%)} \\
    \midrule
    LRU            & 70.6 & 63.2 & 70.0  & 29.4 \\
    LFU            & 70.0 & 66.2 & 70.0  & 30.0 \\
    DCCC\_HASH     & 66.1 & 45.2 & 78.3  & 33.9 \\
    \textbf{DEADLINE\_DCCC} & \textbf{54.5} & \textbf{54.5} & \textbf{87.1} & \textbf{45.5} \\
    \bottomrule
  \end{tabular}
\end{table}

Additionally, the DEADLINE\_DCCC policy produces the following
detailed metric breakdown: edge hit rate 41.4\%, cooperative hit rate
13.1\%, average energy 67.88 (normalised units), total evictions
2\,737, load imbalance 1.76, P95 latency 175.4\,ms, and P99 latency
181.0\,ms.

\subsection{Analysis}

\textbf{Hit rate and deadline satisfaction.}
DEADLINE\_DCCC achieves a 54.5\% overall hit rate and a 54.5\%
deadline satisfaction rate.  In the simulator's implementation, these
two metrics are structurally linked: a cached request is counted as a
hit only when served within the applicable deadline tier (edge hits
within $D_1$; cooperative hits within $D_2$).  Cloud fetches are never
deadline-satisfying.  This design choice means that DEADLINE\_DCCC's
hit rate directly measures \emph{useful} cache hits—those that also
respect latency requirements—rather than raw cache presence.

LRU and LFU achieve higher raw hit rates (70.6\% and 70.0\%) because
they do not enforce deadline filtering: a cache hit is recorded
regardless of the delivery latency.  Under strict deadline-aware
counting, their effective deadline satisfaction rates are 63.2\% and
66.2\% respectively, partly because their simpler eviction policies
tend to keep recently or frequently requested content in cache without
regard to which node can serve it fastest.

\textbf{Cloud fetch and latency.}
DEADLINE\_DCCC's higher cloud fetch rate (45.5\% vs.\ 29.4--33.9\%
for baselines) reflects the conservative deadline enforcement: requests
that baselines serve from cache with latency exceeding $D_2$ are
forwarded to the cloud by DEADLINE\_DCCC.  This increases average
latency (87.1\,ms) compared to LRU/LFU (70.0\,ms) but ensures that
served requests are those that genuinely meet QoS requirements.

\textbf{DCCC\_HASH comparison.}
The DCCC\_HASH policy, which uses the same cooperative caching
infrastructure but selects nodes by consistent hashing, achieves only
a 45.2\% deadline satisfaction rate and 78.3\,ms average latency.
The scoring function in DEADLINE\_DCCC avoids the worst-case latency
of hash routing by routing around loaded or geographically distant
home nodes.

\textbf{Cooperative hit decomposition.}
Of the total hits achieved by DEADLINE\_DCCC, 41.4 percentage points
correspond to edge hits (served by the local/optimal node within
$D_1$) and 13.1 percentage points to cooperative hits (served by a
peer within $D_2$).  This decomposition validates that the two-tier
routing architecture is actively using both tiers, not degenerating to
single-tier operation.

\textbf{Load imbalance.}
The load imbalance metric of 1.76 (defined as the ratio of the
maximum to minimum node load at the end of evaluation) reflects the
heterogeneous capacity configuration: strong nodes naturally attract
more traffic than weak nodes under the scoring function.  A perfectly
balanced load would yield a ratio of 1.0.  Reducing this further
requires increasing the coefficient on the load-squared penalty or
introducing explicit capacity-aware routing, which is left for future
work.

% =========================================================================
\section{Conclusion}
\label{sec:conclusion}

We have presented DEADLINE\_DCCC, a deadline-aware multi-agent deep
reinforcement learning framework for cooperative edge caching under
partial observability.  The key innovations—the POMDP formulation with
independent per-node DQN agents, the Newton-cooling popularity model,
and the six-term deadline-aware scoring function—collectively address
three structural limitations of the ICE base paper: the centralisation
assumption, the static popularity model, and the absence of latency
deadline enforcement.

Our evaluation on a heterogeneous five-node cluster with region-skewed
Zipf traffic demonstrates that DEADLINE\_DCCC achieves a 54.5\%
deadline satisfaction rate with 41.4\% edge hits, 13.1\% cooperative
hits, and an average latency of 87.1\,ms.  While raw hit rates are
lower than LRU/LFU baselines, this reflects a deliberate design choice
to count only deadline-satisfying hits, making the metric directly
meaningful for real-time applications.

Future work will explore (i)~joint training using centralised critics
with decentralised actors (CTDE) to accelerate convergence while
retaining runtime independence, (ii)~integration with EdgeSimPy for
higher-fidelity mobility and energy modelling, (iii)~dynamic coefficient
adaptation through meta-learning to handle workload shifts, and
(iv)~validation on real CDN traces with measured deadline distributions.

% =========================================================================
\section*{Acknowledgement}
The authors thank the Department of Computer Science and Engineering at
Dr.\ B.\ R.\ Ambedkar NIT Jalandhar for providing computational
resources and the DCCC research group for helpful discussions.

% =========================================================================
\begin{thebibliography}{15}

\bibitem{shi2016edge}
W.~Shi, J.~Cao, Q.~Zhang, Y.~Li, and L.~Xu,
``Edge computing: Vision and challenges,''
\emph{IEEE Internet of Things Journal}, vol.~3, no.~5, pp.~637--646, Oct.\ 2016.

\bibitem{ice2022}
A.~Al-Habob, O.~Dobre, H.~Yanikomeroglu, and S.~Muhaidat,
``Towards intelligent adaptive edge caching using deep reinforcement
learning,''
\emph{IEEE Transactions on Vehicular Technology}, vol.~71, no.~8, pp.~8983--8998, Aug.\ 2022.

\bibitem{sadeghi2019optimal}
A.~Sadeghi, F.~Sheikholeslami, and G.~B.~Giannakis,
``Optimal and scalable caching for 5G using reinforcement learning of
space-time popularities,''
\emph{IEEE Journal of Selected Topics in Signal Processing},
vol.~12, no.~1, pp.~180--190, Feb.\ 2019.

\bibitem{zhong2020deep}
C.~Zhong, M.~C.~Gursoy, and S.~Velipasalar,
``A deep reinforcement learning-based framework for content caching,''
in \emph{Proc.\ 52nd Annual Conference on Information Sciences and
Systems (CISS)}, Princeton, NJ, Mar.\ 2020, pp.~1--6.

\bibitem{liu2019decentralised}
J.~Liu, B.~Bai, J.~Zhang, and K.~B.~Letaief,
``Cache placement in fog-RANs: From centralized to distributed
algorithms,''
\emph{IEEE Transactions on Wireless Communications}, vol.~16, no.~11,
pp.~7039--7051, Nov.\ 2017.

\bibitem{traverso2013temporal}
S.~Traverso, M.~Ahmed, M.~Garetto, P.~Giaccone, E.~Leonardi, and
S.~Niccolini,
``Temporal locality in today's content caching: Why it matters and how
to model it,''
\emph{ACM SIGCOMM Computer Communication Review}, vol.~43, no.~5,
pp.~5--12, Oct.\ 2013.

\bibitem{borst2010distributed}
S.~Borst, V.~Gupta, and A.~Walid,
``Distributed caching algorithms for content distribution networks,''
in \emph{Proc.\ IEEE INFOCOM}, San Diego, CA, Mar.\ 2010, pp.~1--9.

\bibitem{golrezaei2013femtocaching}
N.~Golrezaei, K.~Shanmugam, A.~G.~Dimakis, A.~F.~Molisch, and
G.~Caire,
``FemtoCaching: Wireless content delivery through distributed caching
helpers,''
\emph{IEEE Transactions on Information Theory}, vol.~59, no.~12,
pp.~8402--8413, Dec.\ 2013.

\bibitem{zheng2018heterogeneous}
W.~Zheng, L.~Gao, G.~Cao, and H.~Zhang,
``Heterogeneous cellular network with cooperative caching: A matching
and coalition formation game approach,''
\emph{IEEE Transactions on Vehicular Technology}, vol.~67, no.~7,
pp.~6415--6428, Jul.\ 2018.

\bibitem{kaelbling1998planning}
L.~P.~Kaelbling, M.~L.~Littman, and A.~R.~Cassandra,
``Planning and acting in partially observable stochastic domains,''
\emph{Artificial Intelligence}, vol.~101, no.~1--2, pp.~99--134, 1998.

\bibitem{lowe2017multi}
R.~Lowe, Y.~Wu, A.~Tamar, J.~Harb, P.~Abbeel, and I.~Mordatch,
``Multi-agent actor-critic for mixed cooperative-competitive
environments,''
in \emph{Advances in Neural Information Processing Systems (NeurIPS)},
Long Beach, CA, Dec.\ 2017, pp.~6379--6390.

\bibitem{xu2019joint}
J.~Xu, L.~Chen, and P.~Zhou,
``Joint service caching and task offloading for mobile edge computing
in dense networks,''
in \emph{Proc.\ IEEE INFOCOM}, Paris, France, Apr.\ 2019, pp.~207--215.

\bibitem{liu2020deadline}
C.~Liu, M.~Bennis, M.~Debbah, and H.~V.~Poor,
``Dynamic task offloading and resource allocation for ultra-reliable
low-latency edge computing,''
\emph{IEEE Transactions on Communications}, vol.~67, no.~6,
pp.~4132--4150, Jun.\ 2019.

\bibitem{mnih2015human}
V.~Mnih, K.~Kavukcuoglu, D.~Silver, A.~A.~Rusu, J.~Veness,
M.~G.~Bellemare, A.~Graves, M.~Riedmiller, A.~K.~Fidjeland,
G.~Ostrovski, S.~Petersen, C.~Beattie, A.~Sadik, I.~Antonoglou,
H.~King, D.~Kumaran, D.~Wierstra, S.~Legg, and D.~Hassabis,
``Human-level control through deep reinforcement learning,''
\emph{Nature}, vol.~518, pp.~529--533, Feb.\ 2015.

\bibitem{buttazzo2011hard}
G.~C.~Buttazzo,
\emph{Hard Real-Time Computing Systems: Predictable Scheduling
Algorithms and Applications}, 3rd ed.
New York, NY: Springer, 2011.

\end{thebibliography}

\end{document}
